% Generated by scripts/assemble.py; corrections belong in source/corrections.json.
\examhead{166}{四川大学 2026 年数学分析试卷}{158}
% p159-b01; source PDF page 159
\blockanchor{p159-b01}
\question{1}(16 分) 求极限:\par
% p159-b02; source PDF page 159
\blockanchor{p159-b02}
(1) $\lim_{n\to \infty}\left(\frac{1}{\ln(n + \pi) - \ln n} -\frac{n}{\pi}\right).$\par
% p159-b03; source PDF page 159
\blockanchor{p159-b03}
(2) $\lim_{x\to 0^{+}}\frac{\int_0^{x^2}(x - t)(1 - \cos\sqrt{t})\mathrm{d}t}{\int_0^x\tan^4 t\sec t\mathrm{d}t}.$\par
% p159-b04; source PDF page 159
\blockanchor{p159-b04}
\question{2}(12 分) 计算 $I_{n} = \int_{0}^{\pi} \cos^{n} x \, \mathrm{d}x$.\par
% p159-b05; source PDF page 159
\blockanchor{p159-b05}
\question{3}(12 分) 要做一个容器为 $1 \mathrm{~m}^{3}$ 的无盖铝圆桶, 什么样的尺寸才能使用料最省?\par
% p159-b06; source PDF page 159
\par\Needspace{6\baselineskip}
\blockanchor{p159-b06}
\question{4}(12 分) 计算\par
\nopagebreak[4]
% p159-b07; source PDF page 159
\blockanchor{p159-b07}
\displayformula{\int_ {L} (x - y) \mathrm{d} x + z ^ {3} \mathrm{d} z.}
% p159-b08; source PDF page 159
\blockanchor{p159-b08}
其中 $L$ 是 $x^{2} + y^{2} + z^{2} = 1$ 与 $z = c(|c| < 1)$ 的交线, 从 $z$ 轴的正向往负向看去是逆时针方向.\par
% p159-b09; source PDF page 159
\par\Needspace{6\baselineskip}
\blockanchor{p159-b09}
\question{5}(12 分) 计算\par
\nopagebreak[4]
% p159-b10; source PDF page 159
\blockanchor{p159-b10}
\displayformula{\iint_ {\Sigma} \frac {x \mathrm{d} y \mathrm{d} z + y \mathrm{d} z \mathrm{d} x + z \mathrm{d} x \mathrm{d} y}{(x ^ {2} + y ^ {2} + z ^ {2}) ^ {\frac {3}{2}}}.}
% p159-b11; source PDF page 159
\blockanchor{p159-b11}
其中 $\Sigma: 1 - \frac{z}{3} = \frac{(x - 2)^2}{25} + \frac{(y - 1)^2}{16} (z \geq 0)$, 方向取上侧.\par
% p159-b12; source PDF page 159
\blockanchor{p159-b12}
\question{6}(13 分) 求幂级数 $\sum_{n=0}^{\infty} (-1)^{n} \frac{x^{n+3}}{(n+1)(n+3)}$ 的收敛域与和函数.\par
% p159-b13; source PDF page 159
\blockanchor{p159-b13}
\question{7}(13 分) 判断积分 $\int_{1}^{+\infty} x^{3} \sin(x^{q}) \, \mathrm{d}x$ ($q \neq 0$) 的敛散性(包括绝对收敛, 条件收敛和发散).\par
% p159-b14; source PDF page 159
\blockanchor{p159-b14}
\question{8}(12 分) 证明: 方程 $e^x = x^2 + 5x + 1$ 的实根不超过 3 个.\par
% p159-b15; source PDF page 159
\blockanchor{p159-b15}
\question{9}(12 分) 设 $f: [0, +\infty) \to [0, +\infty)$ 连续, 且 $\lim_{x \to +\infty} f(x) \int_0^x f(t) \, \mathrm{d}t = a \in (0, +\infty)$, 证明: $\lim_{x \to +\infty} \sqrt{x} f(x)$ 存在且为正.\par
% p159-b16; source PDF page 159
\blockanchor{p159-b16}
\question{10}(12 分) 设 $f(x)$ 在 $[0, +\infty)$ 上一致连续, 且对任意固定的 $x \in [0, +\infty)$, 当自然数 $n \to +\infty$ 时, 有 $f(x + n) \to 0$. 证明: 函数列 $\{f(x + n)\}_{n=1}^{\infty}$ 在 $[0, 1]$ 上一致收敛于 0.\par
% p159-b17; source PDF page 159
\blockanchor{p159-b17}
\question{11}(14 分) 设 $f(x, y)$ 在 $(x_0, y_0)$ 处可微, 且 $f(x_0, y_0) = 0$. $g(x, y)$ 在 $(x_0, y_0)$ 处连续. 证明: $f(x, y)g(x, y)$ 在 $(x_0, y_0)$ 处可微, 且 $\mathrm{d}(fg)(x_0, y_0) = g(x_0, y_0)\mathrm{d}f(x_0, y_0)$.\par
% p159-b18; source PDF page 159
\blockanchor{p159-b18}
\question{12}(10 分) 设 $f(x)$ 为 [0,1] 上非负且严格单调增加的连续函数. 证明: 对任意的正整数 $n$, 存在唯一的 $x_{n} \in [0,1]$, 使得 $(f(x_{n}))^{n} = \int_{0}^{1}(f(x))^{n}\mathrm{d}x$, 且 $\lim_{n\to \infty}x_n = 1$.\par
